% !TeX program = xelatex
% LINE Messaging API 教學文件：XeLaTeX 原始碼
\documentclass[12pt,a4paper]{report}

\usepackage[margin=2.5cm,headheight=15pt]{geometry}
\usepackage{fontspec}
\usepackage{xeCJK}
\setCJKmainfont{MingLiU}
\usepackage{setspace}
\onehalfspacing
\usepackage{graphicx}
\usepackage{booktabs}
\usepackage{tabularx}
\usepackage{array}
\usepackage{xcolor}
\definecolor{linegreen}{HTML}{06C755}
\definecolor{navy}{HTML}{17324D}
\definecolor{codebg}{HTML}{F4F7FA}
\definecolor{codecomment}{HTML}{527A58}
\definecolor{codeblue}{HTML}{245A9B}
\usepackage{listings}
\lstdefinestyle{pythonstyle}{
  language=Python,
  backgroundcolor=\color{codebg},
  basicstyle=\ttfamily\small,
  keywordstyle=\color{codeblue}\bfseries,
  commentstyle=\color{codecomment},
  stringstyle=\color{red!65!black},
  numbers=left,
  numberstyle=\tiny\color{gray},
  numbersep=9pt,
  stepnumber=1,
  showstringspaces=false,
  breaklines=true,
  breakatwhitespace=true,
  frame=single,
  rulecolor=\color{black!18},
  tabsize=4,
  columns=fullflexible,
  keepspaces=true,
  extendedchars=true,
  captionpos=b
}
\lstdefinestyle{shellstyle}{
  backgroundcolor=\color{codebg}, basicstyle=\ttfamily\small,
  numbers=left, numberstyle=\tiny\color{gray}, numbersep=9pt,
  breaklines=true, frame=single, rulecolor=\color{black!18},
  columns=fullflexible, keepspaces=true, showstringspaces=false,
  captionpos=b
}
\lstdefinestyle{latexstyle}{
  language=[LaTeX]TeX,
  backgroundcolor=\color{codebg}, basicstyle=\ttfamily\scriptsize,
  keywordstyle=\color{codeblue}, commentstyle=\color{codecomment},
  numbers=left, numberstyle=\tiny\color{gray}, numbersep=8pt,
  breaklines=true, breakatwhitespace=false, frame=single,
  rulecolor=\color{black!18}, columns=fullflexible, keepspaces=true,
  showstringspaces=false, extendedchars=true, captionpos=b
}
\usepackage{tikz}
\usetikzlibrary{arrows.meta,positioning,shapes.geometric}
\usepackage{fancyhdr}
\pagestyle{fancy}
\fancyhf{}
\fancyhead[L]{LINE Messaging API 教學文件}
\fancyhead[R]{\nouppercase{\leftmark}}
\fancyfoot[C]{\thepage}
\renewcommand{\headrulewidth}{0.4pt}
\usepackage{hyperref}
\hypersetup{
  colorlinks=true,
  linkcolor=navy,
  urlcolor=blue!65!black,
  citecolor=navy,
  pdftitle={LINE Messaging API 註冊申請與實作教學},
  pdfauthor={簡宏洲}
}
\usepackage{bookmark}
\setlength{\parindent}{2em}
\setlength{\parskip}{3pt}
\renewcommand{\contentsname}{章節目錄}
\renewcommand{\listfigurename}{圖目錄}
\renewcommand{\listtablename}{表目錄}
\renewcommand{\lstlistlistingname}{程式目錄}
\renewcommand{\lstlistingname}{程式碼}
\renewcommand{\chaptermark}[1]{\markboth{第\thechapter 章\quad #1}{}}
\makeatletter
\def\@makechapterhead#1{%
  \vspace*{50\p@}%
  {\parindent \z@ \raggedright \normalfont
    \ifnum \c@secnumdepth >\m@ne
      \huge\bfseries 第\thechapter 章\par\nobreak
      \vskip 20\p@
    \fi
    \interlinepenalty\@M
    \Huge\bfseries #1\par\nobreak
    \vskip 40\p@
  }}
\makeatother

\begin{document}
\begin{titlepage}
  \thispagestyle{empty}
  \begin{tikzpicture}[remember picture,overlay]
    \fill[navy] (current page.north west) rectangle ([yshift=-9.5cm]current page.north east);
    \fill[linegreen] ([yshift=-9.5cm]current page.north west) rectangle ([yshift=-9.64cm]current page.north east);
    \fill[linegreen,opacity=0.13] ([xshift=-1.2cm,yshift=-1.6cm]current page.north east) circle (3.0cm);
    \fill[white,opacity=0.045] ([xshift=-2.0cm,yshift=-2.0cm]current page.north east) circle (4.0cm);
    \fill[navy] (current page.south west) rectangle ([yshift=3.0cm]current page.south east);
    \draw[white,opacity=0.3] ([xshift=2.5cm,yshift=2.35cm]current page.south west) -- ([xshift=-2.5cm,yshift=2.35cm]current page.south east);
    \node[anchor=south west,align=left,text=white] at ([xshift=2.5cm,yshift=0.75cm]current page.south west)
      {\small 課程：專業排版與設計\\[4pt]\large 學號：413631655\quad 作者：簡宏洲};
    \node[anchor=south east,align=right,text=white] at ([xshift=-2.5cm,yshift=0.75cm]current page.south east)
      {\small 作業專題\\[4pt]\large 2026 年 11 月 13 日};
  \end{tikzpicture}
  \centering
  \vspace*{0.4cm}
  \begin{flushleft}
    {\small\bfseries\color{linegreen} 專業排版與設計\quad /\quad 期中專題}\par
    \vspace{0.6cm}
    {\Huge\bfseries\color{white} LINE Messaging API\par}
    \vspace{0.2cm}
    {\huge\bfseries\color{white} 註冊申請與實作教學\par}
    \vspace{0.2cm}
    {\large\color{white!82} 從官方帳號設定到第一個 Python Echo Bot\par}
  \end{flushleft}
  \vspace{2.1cm}
  {\large\bfseries\color{navy} 學習路徑}\par
  \vspace{0.45cm}
  \begin{tikzpicture}[
    node distance=0.62cm,
    flowcard/.style={draw=black!12,rounded corners=5pt,fill=white,minimum width=4.25cm,minimum height=2.8cm,text width=3.75cm,align=center,inner sep=0.15cm},
    flowarrow/.style={-{Stealth[length=2.2mm]},draw=linegreen,thick}
  ]
    \node[flowcard] (account) {\small\bfseries\color{linegreen} STEP 01\\[5pt]\large\bfseries\color{navy} 建立官方帳號\\[3pt]\small\normalfont\color{black!65}準備 Provider 與 Channel};
    \node[flowcard,right=of account] (webhook) {\small\bfseries\color{linegreen} STEP 02\\[5pt]\large\bfseries\color{navy} 設定 Webhook\\[3pt]\small\normalfont\color{black!65}安全接收 LINE 事件};
    \node[flowcard,right=of webhook] (bot) {\small\bfseries\color{linegreen} STEP 03\\[5pt]\large\bfseries\color{navy} Echo Bot 實作\\[3pt]\small\normalfont\color{black!65}驗證後回覆文字訊息};
    \draw[flowarrow] (account)--(webhook);
    \draw[flowarrow] (webhook)--(bot);
  \end{tikzpicture}
  \vfill
\end{titlepage}
\pagenumbering{roman}

\chapter*{摘要}
本文件以繁體中文說明 LINE Messaging API 的申請、設定與入門實作，並以 Python、Flask 及 LINE 官方 Messaging API SDK v3 建立文字回聲機器人。內容涵蓋 LINE Official Account 與 Messaging API Channel 的建立關係、Provider 選擇、Channel Secret 與 Channel access token 的用途、Webhook HTTPS 設定、簽章驗證、ngrok 本機測試，以及常見錯誤的排查方式。本文依據 LINE Developers 官方文件於 2026 年 11 月 13 日查閱；操作介面與平台政策可能調整，實際操作時仍應以官方最新文件為準。

\textbf{關鍵字：} LINE Messaging API、LINE Official Account、Webhook、Python、Flask、XeLaTeX

\tableofcontents
\listoffigures
\listoftables
\lstlistoflistings
\clearpage
\pagenumbering{arabic}

\chapter{LINE Messaging API 概論}
\section{服務概念}
LINE Messaging API 是 LINE Platform 提供給服務開發者的訊息介面。開發者建立 LINE Official Account（官方帳號）並啟用 Messaging API 後，會取得對應的 Channel；伺服器可透過 Webhook 接收使用者事件，再以 Reply 或 Push 等 API 傳送訊息。Channel 是服務與 LINE Platform 之間的設定及認證單位，Provider 則代表管理這些 Channel 的個人或組織。

訊息往返可分成兩條路徑：使用者送出的訊息由 LINE Platform 以 HTTPS POST 傳送至已登記的 Webhook URL；Bot 伺服器驗證請求、處理事件後，使用 Channel access token 呼叫 LINE API 回覆。Webhook 與 API 呼叫是不同方向的連線，因此只開發回覆程式而未設定 Webhook，Bot 仍收不到訊息。

\section{應用情境與訊息種類}
常見應用包含客服問答、訂單或預約通知、校園公告、活動提醒與簡易會員服務。文字訊息可用於 FAQ 或回聲測試；Quick Reply 與 Rich Menu 適合提供固定操作入口；Flex Message 可用卡片呈現商品、行程或訂單摘要。Push Message 可主動通知已符合平台條件的對象，但會受帳號方案、訊息額度與平台政策限制，正式上線前應檢查官方說明。

\begin{table}[htbp]
\centering
\caption{Messaging API 常見功能與用途}
\label{tab:features}
\begin{tabularx}{\textwidth}{>{\raggedright\arraybackslash}p{3.0cm} X X}
\toprule
功能 & 主要用途 & 入門注意事項 \\
\midrule
Webhook & 接收訊息、加入好友等事件 & 必須使用可公開連線的 HTTPS URL，並驗證簽章 \\
Reply Message & 回應觸發事件的使用者 & 使用事件提供的 reply token，應及時處理 \\
Push Message & 對已知對象主動推播 & 需有合法 user ID，並確認方案及訊息額度 \\
Rich Menu／Flex Message & 提供選單或彈性版面 & 先處理互動流程與手機版呈現 \\
\bottomrule
\end{tabularx}
\end{table}

\section{運作流程}
圖 \ref{fig:flow} 概述文字回覆的基本資料流。實際服務還應加入錯誤記錄、逾時處理與重複事件防護；Webhook 可能重送，因此不應假設每個事件只會抵達一次。
\begin{figure}[htbp]
\centering
\begin{tikzpicture}[node distance=0.85cm,>=Stealth,
  box/.style={draw=navy,rounded corners,fill=blue!4,minimum width=3.2cm,minimum height=0.85cm,align=center},
  decision/.style={draw=orange!75!black,diamond,aspect=2,fill=orange!10,align=center,inner sep=2pt}]
  \node[box] (send) {使用者傳送文字};
  \node[box,below=of send] (webhook) {LINE Platform POST Webhook};
  \node[decision,below=of webhook] (verify) {簽章有效？};
  \node[box,left=1.4cm of verify] (reply) {呼叫 Reply API};
  \node[box,right=1.4cm of verify] (reject) {拒絕請求並記錄};
  \draw[->,thick] (send)--(webhook);
  \draw[->,thick] (webhook)--(verify);
  \draw[->,thick] (verify)--node[above,font=\small]{是}(reply);
  \draw[->,thick] (verify)--node[above,font=\small]{否}(reject);
\end{tikzpicture}
\caption{Echo Bot 的 Webhook 與回覆流程}
\label{fig:flow}
\end{figure}

\chapter{申請帳號與建立 Messaging API Channel}
\section{準備項目}
開始前準備可登入的 LINE 帳號或電子郵件、可管理官方帳號的 Business ID，以及可執行 Python 的電腦或伺服器。若只在本機開發，還須準備 ngrok 等 HTTPS 通道工具；正式服務則應使用具有效 TLS 憑證的網域與伺服器。管理者應確認官方帳號的擁有權及 Provider 歸屬，避免將不同組織或服務的 Channel 混在不適當的 Provider 下。

\section{建立官方帳號並啟用 Messaging API}
LINE 官方流程已調整：Messaging API Channel 是透過 LINE Official Account 啟用 Messaging API 時建立，不再由 LINE Developers Console 直接新增 Messaging API Channel。建議依下列順序操作：
\begin{enumerate}
  \item 開啟 LINE Official Account 的申請頁面，使用 Business ID 註冊或登入。
  \item 填寫官方帳號名稱、產業類別及必要資料，完成官方帳號建立。
  \item 前往 LINE Official Account Manager，選擇剛建立的帳號並確認基本資料。
  \item 在帳號設定中啟用 Messaging API。首次使用 Developers Console 時，依畫面登記開發者姓名與電子郵件。
  \item 選擇管理此官方帳號的 Provider。Provider 指派後通常不能任意更換，應先確認由個人或組織管理。
  \item 登入 LINE Developers Console，開啟剛選定的 Provider，確認系統已建立 Messaging API Channel。
\end{enumerate}

\begin{figure}[htbp]
\centering
\begin{tikzpicture}[node distance=0.65cm,
  flowstep/.style={draw=linegreen!80!black,rounded corners,fill=linegreen!7,minimum width=11cm,minimum height=0.68cm,align=center},
  arr/.style={-{Stealth},thick,linegreen!80!black}]
  \node[flowstep] (a) {Business ID 註冊／登入};
  \node[flowstep,below=of a] (b) {建立 LINE Official Account};
  \node[flowstep,below=of b] (c) {於 Official Account Manager 啟用 Messaging API};
  \node[flowstep,below=of c] (d) {選擇 Provider，系統建立 Messaging API Channel};
  \node[flowstep,below=of d] (e) {至 Developers Console 管理 Channel 設定};
  \draw[arr] (a)--(b); \draw[arr] (b)--(c); \draw[arr] (c)--(d); \draw[arr] (d)--(e);
\end{tikzpicture}
\caption{建立 Messaging API Channel 的現行流程}
\label{fig:registration}
\end{figure}

\section{Provider 與 Channel 的管理原則}
一個 Provider 可管理多個 Channel；Channel 各自代表一種平台服務及其設定。Provider 的選擇會影響管理權限及使用者識別範圍，且官方流程可能限制日後移轉，因此個人作業可用個人 Provider，團隊或企業服務則應由組織授權的管理者建立。不要把 Channel Secret、Access Token 或任何登入驗證資訊貼在公開文件、截圖或程式碼庫。

\chapter{Channel 憑證與安全設定}
\section{Channel Secret 與 Access Token}
Channel Secret 用於驗證 Webhook 請求的 HMAC 簽章；Channel access token 則用於授權伺服器呼叫 Messaging API。兩者用途不同，不能互相替代。Secret 由 Channel 的 Basic settings 頁籤查看；Access Token 在 Messaging API 頁籤依介面選擇並簽發。官方提供多種 Token 型態，本文建議依官方最新文件優先評估具自訂到期時間的 Channel access token v2.1，並建立輪替與撤銷程序。

\begin{table}[htbp]
\centering
\caption{Channel 憑證用途與保護方式}
\label{tab:credentials}
\begin{tabularx}{\textwidth}{>{\raggedright\arraybackslash}p{3.2cm} X X}
\toprule
憑證 & 用途 & 安全措施 \\
\midrule
Channel Secret & 驗證 Webhook 的 X-Line-Signature & 存於環境變數或秘密管理服務；不得寫入日誌 \\
Channel access token & 呼叫 Reply、Push 等 Messaging API & 限定權限及期限，定期輪替，外洩時立即撤銷 \\
Reply token & 回應單一 Webhook 事件 & 事件專用且有使用限制；收到事件後盡快回覆 \\
\bottomrule
\end{tabularx}
\end{table}

\section{以環境變數保存秘密}
本範例從作業系統環境變數讀取憑證，不將真實值寫入原始碼。開發環境可使用未納入版本控制的 \texttt{.env} 檔，並設定 \texttt{.gitignore} 忽略該檔；部署環境則使用平台提供的 Secret／Environment Variables。若憑證曾被提交，即使後續刪除檔案，Git 歷史中仍可能保留，應立即撤銷或輪替憑證，不能只靠新增忽略規則處理。

\begin{lstlisting}[style=shellstyle,caption={環境變數範例（請填入自己的憑證）},label={lst:env}]
# .env 範例；此檔不得提交至公開 Git 儲存庫
LINE_CHANNEL_SECRET=請填入Channel的Secret
LINE_CHANNEL_ACCESS_TOKEN=請填入Channel的AccessToken
\end{lstlisting}

\chapter{Webhook 設定與測試}
\section{設定 Webhook URL}
在 LINE Developers Console 開啟 Provider 與 Messaging API Channel，進入 Messaging API 頁籤，設定 Webhook URL 為伺服器的公開 HTTPS 網址，例如 \url{https://example.ngrok-free.app/callback}。Webhook URL 必須能由 LINE Platform 連線，使用受信任憑證的 HTTPS；完成後按 Verify 測試連線，再啟用 Use webhook。若測試失敗，先確認伺服器已啟動、路徑正確、TLS 憑證有效，以及防火牆沒有阻擋請求。

\section{本機開發搭配 ngrok}
Flask 本機服務通常只綁定開發電腦，LINE Platform 無法直接連入。ngrok 可建立從公開 HTTPS 網址轉送至本機埠號的通道。以下操作以開發測試為例：
\begin{enumerate}
  \item 安裝套件：Flask、LINE Messaging API Python SDK v3 與 python-dotenv。
  \item 在本機啟動 Flask，確認 \url{http://127.0.0.1:8000/health} 可回傳成功狀態。
  \item 啟動 ngrok，將外部 HTTPS 流量轉送至本機 8000 埠。
  \item 複製 ngrok 顯示的 HTTPS 網址，在尾端加入 \texttt{/callback} 並填入 Webhook URL。
  \item 按 Verify；成功後開啟 Use webhook，將官方帳號加入好友並傳送文字測試。
\end{enumerate}

\begin{lstlisting}[style=shellstyle,caption={安裝及啟動本機測試服務},label={lst:commands}]
python -m venv .venv
# Windows PowerShell
.\.venv\Scripts\Activate.ps1
python -m pip install Flask line-bot-sdk python-dotenv
python app.py
# 另開終端機執行 ngrok http 8000
\end{lstlisting}

免費隧道網址可能在重新啟動後改變，網址變更時須更新 Console 的 Webhook URL。ngrok 僅供開發測試，不應視為正式部署方案；正式環境需使用穩定網域、TLS、監控及適當的網路安全設定。

\section{測試與事件處理注意事項}
Verify 測試成功僅代表端點可連線，並不保證所有事件處理邏輯正確。應從實際 LINE 帳號傳送文字，檢查伺服器收到的事件與回覆結果。Webhook 請求必須先以原始 request body 和 Channel Secret 驗證 X-Line-Signature，再解析或處理事件；不得只根據來源 IP 或自訂欄位判定請求可信。若啟用 Webhook redelivery，同一事件可能重複送達，可用 \texttt{webhookEventId} 實作冪等處理；不支援的事件類型也應安全略過。

\chapter{Python Flask Echo Bot 實作}
\section{專案結構與相依套件}
下列程式以 Flask 提供 \texttt{/callback} Webhook，使用 LINE 官方 Python SDK v3 的 \texttt{WebhookHandler} 驗證簽章並處理文字事件。範例只回傳收到的文字，適合作為連線測試起點。專案可規劃成 \texttt{app.py}、\texttt{.env} 與 \texttt{.gitignore}；本文件提供完整程式，憑證值需由操作者自行設定。

\begin{lstlisting}[style=pythonstyle,caption={Python Flask LINE Echo Bot（app.py）},label={lst:app}]
import os

from dotenv import load_dotenv
from flask import Flask, abort, jsonify, request
from linebot.v3 import WebhookHandler
from linebot.v3.exceptions import InvalidSignatureError
from linebot.v3.messaging import (
    ApiClient,
    Configuration,
    MessagingApi,
    ReplyMessageRequest,
    TextMessage,
)
from linebot.v3.webhooks import MessageEvent, TextMessageContent

load_dotenv()

CHANNEL_SECRET = os.environ.get("LINE_CHANNEL_SECRET")
CHANNEL_ACCESS_TOKEN = os.environ.get("LINE_CHANNEL_ACCESS_TOKEN")
if not CHANNEL_SECRET or not CHANNEL_ACCESS_TOKEN:
    raise RuntimeError("請設定 LINE_CHANNEL_SECRET 與 LINE_CHANNEL_ACCESS_TOKEN")

app = Flask(__name__)
handler = WebhookHandler(CHANNEL_SECRET)
configuration = Configuration(access_token=CHANNEL_ACCESS_TOKEN)


@app.get("/health")
def health():
    return jsonify(status="ok"), 200


@app.post("/callback")
def callback():
    signature = request.headers.get("X-Line-Signature", "")
    body = request.get_data(as_text=True)
    try:
        # SDK 會先驗證簽章，再呼叫已註冊的事件處理函式。
        handler.handle(body, signature)
    except InvalidSignatureError:
        app.logger.warning("Webhook signature verification failed")
        abort(400)
    return "OK", 200


@handler.add(MessageEvent, message=TextMessageContent)
def handle_text_message(event):
    reply_text = f"收到：{event.message.text}"
    reply_request = ReplyMessageRequest(
        reply_token=event.reply_token,
        messages=[TextMessage(text=reply_text)],
    )
    with ApiClient(configuration) as api_client:
        messaging_api = MessagingApi(api_client)
        messaging_api.reply_message(reply_request)


if __name__ == "__main__":
    # 開發測試使用；正式環境請使用 WSGI 伺服器並關閉 debug。
    app.run(host="127.0.0.1", port=8000, debug=False)
\end{lstlisting}

\section{程式流程說明}
\begin{enumerate}
  \item \texttt{load\_dotenv()} 載入本機環境變數；部署時也可由作業系統直接注入變數。
  \item \texttt{WebhookHandler} 使用 Channel Secret 驗證 LINE 傳來的簽章。若簽章無效，程式回傳 HTTP 400，不處理事件。
  \item \texttt{handler.add} 僅註冊文字訊息事件，避免把貼圖、圖片等非文字內容當成文字讀取。
  \item 回覆時使用事件中的 \texttt{reply\_token}，以 SDK 建立 \texttt{ReplyMessageRequest} 並呼叫 Messaging API。
  \item \texttt{/health} 提供簡單健康檢查；它不取代 Webhook Verify，也不應回傳任何 Token 或 Secret。
\end{enumerate}

\section{安全與上線前檢查}
此程式是教學起點，尚未涵蓋背景佇列、資料庫、速率限制及完整觀測性。正式服務可快速驗證簽章後回覆 2xx，再將事件交由佇列處理，降低平台重送風險；同時以 \texttt{webhookEventId} 避免重複執行副作用。部署時使用 Gunicorn 等 WSGI 伺服器或平台建議的執行方式，不要使用 Flask 開發伺服器對外服務。錯誤日誌可記錄事件類型及追蹤編號，但不得記錄完整憑證、簽章或不必要的個人資料。

\begin{table}[htbp]
\centering
\caption{上線前測試檢查表}
\label{tab:checklist}
\begin{tabularx}{\textwidth}{>{\centering\arraybackslash}p{1.3cm} X X}
\toprule
完成 & 檢查項目 & 預期結果 \\
\midrule
□ & 環境變數已設定且不在 Git 追蹤中 & 程式可啟動，秘密不出現在儲存庫 \\
□ & Webhook 使用有效 HTTPS URL & Console Verify 顯示成功 \\
□ & Use webhook 已啟用 & 傳送文字後伺服器收到事件 \\
□ & 無效簽章測試 & 端點拒絕請求，不呼叫回覆 API \\
□ & 正常文字訊息測試 & Bot 以 Reply API 回覆一次 \\
\bottomrule
\end{tabularx}
\end{table}

\chapter{故障排除與延伸應用}
\section{常見問題}
\begin{description}
  \item[Verify 失敗] 確認公開網址可連線、使用 HTTPS、路徑是 \texttt{/callback}，且 Flask 路由接受 POST。
  \item[收到 400] 通常是簽章驗證失敗。確認使用正確 Channel Secret，且在驗證之前沒有改寫 request body。
  \item[回覆 API 失敗] 檢查 Access Token 是否屬於同一個 Channel、是否有效，以及 reply token 是否逾時或已使用。
  \item[傳訊息沒有 Webhook] 確認 Use webhook 已開啟，官方帳號與 Messaging API Channel 關聯正確，並檢查伺服器日誌。
  \item[Bot 出現重複回覆] 檢查官方帳號自動回應設定，並確認程式是否因 Webhook 重送而重複處理事件。
\end{description}

\section{可延伸功能}
完成 Echo Bot 後，可依需求擴充指令路由、FAQ 查詢、Rich Menu、Quick Reply 或 Flex Message。推播通知需先確認使用者同意及帳號方案限制；大型服務可導入資料庫、佇列、重試及事件去重機制。任何擴充都應遵守最小化蒐集、目的限定及保存期限等個資保護原則，並提供適當的使用者告知與退出方式。

\chapter{結論}
本文件從 LINE Official Account 建立開始，說明了 Messaging API Channel 的現行建立方式、Provider 選擇、憑證管理、Webhook HTTPS 設定與 Python Flask 回聲機器人。完成入門測試的關鍵是：正確的 Channel 憑證、可公開連線的 HTTPS Webhook、先驗證簽章再處理事件，以及使用有效 reply token 呼叫 Reply API。進一步部署時，應加入安全的秘密管理、事件冪等、監控與正式 WSGI 伺服器，並持續依 LINE 官方文件更新設定。

\section*{官方參考資料}
\begin{enumerate}
  \item LINE Developers，\href{https://developers.line.biz/en/docs/messaging-api/getting-started/}{Get started with the Messaging API}。
  \item LINE Developers，\href{https://developers.line.biz/en/docs/messaging-api/building-bot/}{Build a bot}。
  \item LINE Developers，\href{https://developers.line.biz/en/docs/messaging-api/receiving-messages/}{Receive messages (webhook)}。
  \item LINE Developers，\href{https://developers.line.biz/en/docs/messaging-api/verify-webhook-signature/}{Verify webhook signature}。
  \item LINE Developers，\href{https://github.com/line/line-bot-sdk-python}{LINE Messaging API SDK for Python}。
\end{enumerate}

\chapter{本文件 LaTeX 原始碼}
以下以 \texttt{listings} 套件直接載入本文件原始碼，設定標題並列入程式目錄。由於此章展示整份文件，頁數會隨原始碼長度增加；這也使作業原始碼與 PDF 內容保持一致。
\lstinputlisting[style=latexstyle,caption={本文件 LaTeX 原始碼（main.tex）},label={lst:latex-source}]{main.tex}

\end{document}
